%!TEX program = xelatex
\documentclass[twocolumn]{article}

\usepackage[tikz]{multicolrule}
\setlength{\columnsep}{24pt}
\SetMCRule{line-style=dots, width=0.6pt}

\usepackage[a4paper,
            left=20mm, right=20mm,
            top=20mm, bottom=20mm]{geometry}

\usepackage{kotex}
\usepackage{fontspec}

\setmainhangulfont{Nanum Myeongjo}
\setsanshangulfont{Nanum Gothic}

\usepackage{amsmath}
\usepackage{amssymb}
\usepackage{amsthm}
\usepackage{graphicx}

\renewcommand{\proofname}{\textsf{\textbf{증명}}\nopunct}
\renewcommand{\qedsymbol}{$\blacksquare$}

\usepackage{setspace}

\begin{document}
\title{2차 함수와 2차 방정식}
\author{}
\date{}

\maketitle

\setstretch{1.25}
\section{2차 함수}
2차 함수는 다항식으로 이뤄진 함수 중에서도 최고차항의 차수가 2인 함수로
일반적으로 다음과 같은 형태로 표현된다.
\[
f(x) := ax^2 + bx + c\ (a, b, c \in \mathbb{R}, a \neq 0)
\]
여기서 \(a\)의 값에 따라 함수의 그래프 형태가 크게 두 가지의 형태로 나뉜다.
\begin{itemize}
    \item \(a > 0\)인 경우: 그래프는 아래로 볼록한 포물선 형태를 가진다.
    \item \(a < 0\)인 경우: 그래프는 위로 볼록한 포물선 형태를 가진다.
\end{itemize}
이는 2차 함수의 그래프가 포물선 형태를 갖고 있음을 숙지하고 있으면 그 방향은
충분히 크고 작은 값을 $x$에 대입하여 확인할 수 있다.

\subsection{2차 함수의 표현식}
2차 함수는 식을 방식에 따라 변형함으로써 함수의 그래프 개형을 더욱 파악하기 쉽다.
\subsubsection{완전제곱식}
완전제곱식이란 아래와 같은 꼴의 식을 말한다.
\[
(ax + b)^2
\]
2차 함수를 완전제곱식으로 묶으면 다음과 같은 꼴로 표현할 수 있다.
\[
f(x) := (ax + b)^2 + m
\]
위 식에서 완전제곱식이 0이 되도록 하는 $x$의 값은 $-\frac{b}{a}$이며
이 때의 함수값은 $m$이 되고 이 $x$값을 기준으로 좌우 대칭임을 알 수 있다.

\begin{proof}
\[
f(-\frac{b}{a} + k) = f(-\frac{b}{a} - k) = k^2 + m\ (\forall k \in \mathbb{R})
\]
\end{proof}

\subsubsection{인수분해식}
인수분해란 다항식을 더 낮은 차수를 갖는 다항식들의 곱으로 표현하는 것을 말한다.
따라서 2차 함수는 최고차수가 2이기 때문에 인수분해를 하면 두 1차식의 곱으로 나타난 형태를 갖는다.
예로 아래와 같이 두 1차식이 곱해진 식을 보자.
\[
(a_1x + b_1)(a_2x + b_2)
\]
이를 전개하면 다음과 같다.
\[
a_1a_2x^2 + (a_1b_2 + a_2b_1)x + b_1b_2
\]
이를 역으로 다시 두 1차식으로 곱해진 식으로 변환하는 것이 2차식의 인수분해가 된다.
다시 2차 함수로 돌아가 위의 두 식을 아래와 같이 함수로 써보자.
\begin{align*}
f(x) & :=a_1a_2x^2 + (a_1b_2 + a_2b_1)x + b_1b_2 \\
&\,= (a_1x + b_1)(a_2x + b_2)
\end{align*}
위 함수에서 우리는 아래와 같음을 알 수 있다.
\[
f(-\frac{b_1}{a_1}) = f(-\frac{b_2}{a_2}) = 0
\]
또한 우리는 이를 통해 위의 함수를 그래프로 표현했을 때,
$x=-\frac{b_1}{a_1}$인 지점과 $x=-\frac{b_2}{a_2}$인 지점에서
$x$축 $(y=0)$과 만난다는 사실을 알 수 있다.
\section{2차 방정식}
2차 방정식은 위에서의 2차 함수가 $x$축과 만나는 지점을 찾는 문제로 다음과 같은 꼴을 갖는다.
\[
ax^2 + bx + c = 0
\]
이를 푸는 방법은 대표적으로 두 가지의 방식이 있는데 그 중 하나는 위에서 설명한 인수분해와
다른 하나로 근의 공식을 이용하는 방법이 있다.
\subsection{근의 공식}
근의 공식은 위와 같은 일반식을 변형해 얻어낼 수 있는데 그 과정은 다음과 같다.
\[
    ax^2 + bx + c = a\left(x^2 + \frac{b}{a}x\right) + c = 0
\]
위에서 우측 두 변에서 양변에 $-c$를 더하고 $a$로 나누면 다음과 같다.
\[
    x^2 + \frac{b}{a}x = - \frac{c}{a}
\]
위에서 양변에 $\left(\frac{b}{2a}\right)^2$을 더하면 좌변은 아래와 같이 완전제곱식이 된다.
\[
    \left(x + \frac{b}{2a}\right)^2 = - \frac{c}{a} + \left(\frac{b}{2a}\right)^2
\]
위에서 양변에 제곱근을 취하고 양변에 $-\frac{b}{2a}$를 더하면 아래와 같은 근의 공식을 얻어낼 수 있다.
\[
    x = \frac{-b \pm \sqrt{b^2 - 4ac}}{2a}
\]


\end{document}